All 27 models were evaluated on all 1,666 annotated external images. The scalar baseline has signed Spearman correlation -0.142 and reader-pair agreement 42.5\%. The scalar conditional mild-blur control achieves 0.348 and 74.5\%, respectively, whereas its unconditional counterpart achieves -0.075 and 49.3\%. This supports a contribution from quality conditioning in this scalar mild-blur comparison, without demonstrating population-level generalization.

The original ordinal strong-blur constraint achieves correlation 0.426 and pair agreement 73.7\%, compared with -0.222 and 37.6\% without the constraint. However, the matched unconditional strong-blur arm achieves 0.391 and 72.8\%. The conditional-minus-unconditional intervals include zero for both endpoints (Table~\ref{tab:paired}); this arm does not isolate a reliable additional benefit from the quality gate.

The ordinal mild-blur arm has correlation -0.097 and pair agreement 48.4\%, despite eliminating score increases on the illustrated synthetic blur probes. ResNet-18 conditioning changes correlation from -0.603 to -0.470, but agreement remains poor (11.8\%). Synthetic constraint satisfaction therefore does not predict successful expert-rated transfer across heads and backbones.

Variation is substantial across training seeds and phantom anatomies (Table~\ref{tab:anatomy} and Figure~\ref{fig:external}). In particular, pooled gains do not establish stable pelvic performance. Paired source-slice intervals are descriptive summaries conditional on the two phantoms and the three fixed seeds; they do not capture uncertainty across patients, sites, physical phantoms or training-seed populations. Adaptive follow-up controls and multiple examined contrasts further limit confirmatory interpretation.
